\section{Conclusiones}

La Actividad 4 convirtió la estructura abstracta de la entrega anterior en una experiencia
visible y recorrible. Al mantener una identidad común, los cinco flujos mostraron que cada
prototipo puede distinguirse por su tarea, información e interacción sin alterar la apariencia
general del producto.

El sistema alcanzó 30 pantallas conectadas, 12 conjuntos de componentes y 75 variantes. La
revisión final eliminó recortes de texto, fuentes faltantes, destinos inválidos y objetivos
interactivos menores de 44 píxeles. Con ello se obtuvo una base consistente para recibir
comentarios sobre la estética y evaluar las tareas con participantes, sin que defectos básicos
de construcción distorsionen la experiencia.

\uxsection{Referencias}

\begin{uxreferencias}
Allen, J., y Chudley, J. (2012). \textit{Smashing UX design: Foundations for designing online
user experiences}. Wiley.

Hartson, R., y Pyla, P. S. (2019). \textit{The UX book: Agile UX design for a quality user
experience} (2.\textsuperscript{a} ed.). Morgan Kaufmann.

Ramírez Mejía, A. I. (2023). \textit{Grocery shopping app: A guided UX case study. Part 4:
Prototyping and Style Guide} [Diapositivas]. TC4032, Instituto Tecnológico y de Estudios
Superiores de Monterrey.
\end{uxreferencias}
